\paragrafos{Είδη γωνιών και κάθετες ευθείες}{Β.1.6, «Είδη γωνιών – Κάθετες ευθείες»}

\begin{anagke}
Η Α3 πρέπει να στήσει τον φράχτη ώστε ο κήπος να έχει «σωστές γωνίες» στις γωνιές του. «Σωστή γωνία σημαίνει τετράγωνη, όπως στο θρανίο», λέει ο Νίκος. Η Αγγελική το λέει επίσημα: ορθή γωνία. Αλλά πώς φτιάχνουν ορθή γωνία στο έδαφος, χωρίς μοιρογνωμόνιο; Ο παππούς δείχνει ένα κόλπο που ξέρουν οι χτίστες από την αρχαιότητα: ένα σχοινί με $12$ ίσα διαστήματα, με κόμπους. Το διπλώνει σε τρίγωνο με πλευρές $3$, $4$ και $5$ διαστημάτων: «Η γωνία ανάμεσα στο $3$ και στο $4$ είναι ορθή».

Η Πολυξένη το καταγράφει και ρωτάει για όλες τις άλλες γωνίες που βλέπει: οξείες, αμβλείες, όπως στο φτερό του αεροπλάνου και στη στέγη. Ο Βασίλης θέλει να ταξινομήσει τις γωνίες του κήπου: «Αυτή που ανοίγει λίγο», «αυτή που ανοίγει πολύ», «η ευθεία». Η κυρία Αλεξάνδρα του δίνει ονόματα. Και όταν δύο μονοπάτια τέμνονται σε ορθή γωνία, τα λέμε κάθετα.

Σε αυτή την παράγραφο ταξινομούμε τις γωνίες ανάλογα με το μέτρο τους και ορίζουμε την καθετότητα, μια από τις πιο χρήσιμες έννοιες στη Γεωμετρία.
\end{anagke}

\begin{lexi}{ορθή, οξεία, αμβλεία, κάθετος}{ορθή, οξεία, αμβλεία}
\lexiword{Ορθή} σημαίνει «σωστή, όρθια» (από το «ορθός»): η γωνία ενός όρθιου τοίχου με το δάπεδο. \lexiword{Οξεία} είναι η «αιχμηρή», σαν την οξεία μύτη. \lexiword{Αμβλεία} είναι η «αμβλυμένη», η στρογγυλεμένη, που ανοίγει περισσότερο από την ορθή. \lexiword{Κάθετη} είναι η ευθεία που πέφτει «κάθετα», σαν το νήμα της στάθμης.
\end{lexi}

\begin{metafrasi}
Γωνία μικρότερη από την ορθή. & οξεία: $0^\circ<\omega<90^\circ$\\
Γωνία $90^\circ$. & ορθή\\
Γωνία ανάμεσα στην ορθή και την ευθεία. & αμβλεία: $90^\circ<\omega<180^\circ$\\
Γωνία $180^\circ$. & ευθεία γωνία\\
Δύο ευθείες που τέμνονται σε ορθή γωνία. & $\varepsilon_1\perp\varepsilon_2$\\
\end{metafrasi}

\tmima{Η θεωρία}

\begin{orismos}{είδη γωνιών}
Μια γωνία $\omega$ λέγεται:
\begin{itemize}[leftmargin=1.2em,itemsep=0.3ex,topsep=2pt]
\item \textbf{οξεία} αν $0^\circ<\omega<90^\circ$,
\item \textbf{ορθή} αν $\omega=90^\circ$,
\item \textbf{αμβλεία} αν $90^\circ<\omega<180^\circ$,
\item \textbf{ευθεία} αν $\omega=180^\circ$ (οι πλευρές της είναι αντικείμενες ημιευθείες),
\item \textbf{μη κυρτή} αν $180^\circ<\omega<360^\circ$ και \textbf{πλήρης} αν $\omega=360^\circ$,
\item \textbf{μηδενική} αν $\omega=0^\circ$ (οι πλευρές συμπίπτουν).
\end{itemize}
\end{orismos}

\begin{center}
\begin{tikzpicture}[scale=0.8]
\foreach \a/\n/\x in {40/οξεία/0,90/ορθή/3.6,130/αμβλεία/7.2}{
\begin{scope}[xshift=\x cm]
\draw[thick] (2.2,0)--(0,0)--({2.2*cos(\a)},{2.2*sin(\a)});
\draw[gblue] (0.6,0) arc(0:\a:0.6);
\node at (1.1,-0.7) {\n};
\end{scope}}
\draw[gblue] (3.6+0.25,0)--(3.6+0.25,0.25)--(3.6,0.25);
\end{tikzpicture}
\end{center}

\tmimasmall{Κάθετες ευθείες}

Δύο ευθείες που τέμνονται σχηματίζουν τέσσερις γωνίες. Αν μία από αυτές είναι ορθή, τότε είναι ορθές και οι τέσσερις. Τότε οι ευθείες λέγονται \textbf{κάθετες} και γράφουμε $\varepsilon_1\perp\varepsilon_2$.

\begin{idiotita}{ύπαρξη και μοναδικότητα κάθετης}
Από ένα σημείο του επιπέδου διέρχεται \textbf{μία και μοναδική} ευθεία κάθετη προς μια δοσμένη ευθεία $(\varepsilon)$, είτε το σημείο ανήκει στην ευθεία είτε όχι.
\end{idiotita}

Τις κάθετες τις σχεδιάζουμε με \textbf{γνώμονα} (τρίγωνο με ορθή γωνία) και χάρακα.

\begin{prosoxi}
\begin{itemize}[leftmargin=1.2em,itemsep=0.3ex,topsep=2pt]
\item Η ορθή γωνία είναι ακριβώς $90^\circ$. Η $91^\circ$ είναι αμβλεία και η $89^\circ$ οξεία.
\item Οι κάθετες δεν χρειάζεται να είναι «οριζόντιες και κατακόρυφες»: δύο ευθείες που βρίσκονται σε οποιαδήποτε θέση είναι κάθετες αν σχηματίζουν ορθή γωνία.
\item Η ευθεία γωνία δεν είναι «ευθεία γραμμή» αλλά γωνία $180^\circ$ με αντικείμενες πλευρές.
\item Από σημείο εκτός ευθείας διέρχεται \emph{μία} κάθετη, αλλά άπειρες πλάγιες.
\end{itemize}
\end{prosoxi}

\tmima{Λυμένα παραδείγματα}

\paradeigma{Ταξινομώ γωνίες}
\tixero Οι γωνίες $37^\circ$, $90^\circ$, $123^\circ$, $180^\circ$, $250^\circ$, $360^\circ$.\par
\tithelo Να πω το είδος κάθε μίας.\par
\possindeo Συγκρίνω με τα $90^\circ$, $180^\circ$, $360^\circ$.\par
\apantisi $37^\circ$ οξεία, $90^\circ$ ορθή, $123^\circ$ αμβλεία, $180^\circ$ ευθεία, $250^\circ$ μη κυρτή, $360^\circ$ πλήρης.\par
\elegxos Η $250^\circ$ είναι ανάμεσα στο $180^\circ$ και στο $360^\circ$.

\paradeigma{Το ρολόι}
\tixero Οι δείκτες στις $2{:}00$, $3{:}00$, $4{:}00$.\par
\tithelo Το είδος της γωνίας.\par
\possindeo Κάθε ώρα είναι $30^\circ$. Στις $2{:}00$: $60^\circ$ οξεία. Στις $3{:}00$: $90^\circ$ ορθή. Στις $4{:}00$: $120^\circ$ αμβλεία.\par
\apantisi Οξεία, ορθή, αμβλεία.\par
\elegxos Στις $6{:}00$ θα ήταν ευθεία γωνία.

\paradeigma{Κάθετες ευθείες}
\tixero Δύο ευθείες που τέμνονται σχηματίζοντας γωνία $90^\circ$.\par
\tithelo Οι υπόλοιπες γωνίες.\par
\possindeo Οι διαδοχικές γωνίες ανήκουν σε ευθεία γωνία $180^\circ$: $180^\circ-90^\circ=90^\circ$. Έτσι και οι τέσσερις είναι ορθές.\par
\apantisi $90^\circ$ η καθεμία.\par
\elegxos Το άθροισμα και των τεσσάρων είναι $360^\circ$.

\paradeigma{Η ορθή γωνία του παππού}
\tixero Σχοινί με $12$ ίσα διαστήματα, διπλωμένο σε τρίγωνο με πλευρές $3$, $4$, $5$.\par
\tithelo Γιατί σχηματίζεται ορθή γωνία.\par
\possindeo Το τρίγωνο με πλευρές $3$, $4$, $5$ έχει ορθή γωνία ανάμεσα στις πλευρές $3$ και $4$ (το αποδεικνύουμε αργότερα). Οι αρχαίοι αρπεδονάπτες το χρησιμοποιούσαν για τις γωνίες των κτιρίων και των αγρών.\par
\apantisi Γιατί $3^2+4^2=5^2$ ($9+16=25$).\par
\elegxos Δοκίμασε με γνώμονα.

\tmima{Ασκήσεις}

\zoni{Το κατάλαβα}{είδη γωνιών και κάθετες}

\ex Ταξινόμησε τις γωνίες $37^\circ$, $90^\circ$, $123^\circ$, $180^\circ$, $250^\circ$, $360^\circ$, $0^\circ$, $89^\circ$.

\ex Σχεδίασε με μοιρογνωμόνιο μια οξεία, μια ορθή και μια αμβλεία γωνία και μέτρησε τις διχοτόμους τους.

\ex Βρες το είδος της γωνίας των δεικτών του ρολογιού στις $2{:}00$, $3{:}00$, $4{:}00$, $6{:}00$, $7{:}00$ (πάρε τη μικρότερη).

\ex Δύο ευθείες τέμνονται και μια γωνία είναι $90^\circ$. Πόσες είναι οι άλλες γωνίες;

\ex Με γνώμονα σχεδίασε μια ευθεία $(\varepsilon)$, ένα σημείο $\mathrm{Α}$ πάνω της και ένα σημείο $\mathrm{Β}$ εκτός. Φέρε από το $\mathrm{Α}$ και από το $\mathrm{Β}$ τις κάθετες προς την $(\varepsilon)$.

\ex Πόσες κάθετες διέρχονται από ένα σημείο προς μια ευθεία; Και πόσες πλάγιες;

\ex Σωστό ή λάθος; α) Μια γωνία $91^\circ$ είναι αμβλεία. \quad β) Η ευθεία γωνία είναι $180^\circ$. \quad γ) Δύο κάθετες ευθείες σχηματίζουν γωνίες που δεν είναι όλες ίσες. \quad δ) Η $0^\circ$ είναι οξεία. \quad ε) Η $250^\circ$ είναι κυρτή.

\zoni{Το χρησιμοποιώ}{γωνιές, στέγες, αναβάτες}

\ex \textbf{Ο φράχτης.} Στις γωνιές του κήπου ο παππούς θέλει ορθές γωνίες. Πώς θα το ελέγξεις με ένα φύλλο χαρτί $\mathrm{Α4}$;

\ex \textbf{Το σχοινί του παππού.} Ένα σχοινί με $12$ ίσα διαστήματα διπλώνεται σε τρίγωνο με πλευρές $3$, $4$, $5$. Τι γωνία σχηματίζεται και πού;

\ex \textbf{Τα δωμάτια.} Βρες στην τάξη τρεις ορθές, δύο οξείες και μία αμβλεία γωνία.

\ex \textbf{Η στέγη.} Μια στέγη έχει γωνία $110^\circ$ στην κορυφή. Τι είδους είναι; Αν η στέγη ήταν ορθή γωνία πόσες μοίρες θα ήταν;

\ex \textbf{Ο ανεμοδείκτης.} Ένας ανεμοδείκτης στρέφεται από τη διεύθυνση Βορράς προς Ανατολή. Τι γωνία διέγραψε; Και προς Νότο;

\zoni{Το σκέφτομαι}{όταν μία ορθή φέρνει τρεις}

\ex Δείξε ότι αν μία από τις γωνίες που σχηματίζουν δύο τεμνόμενες ευθείες είναι ορθή, είναι ορθές και οι άλλες τρεις.

\ex Είναι δυνατόν να έχει ένα τρίγωνο δύο ορθές γωνίες; Δοκίμασε να σχεδιάσεις.

\ex Γιατί ο γνώμονας φτιάχνει κάθετες; Πώς θα έλεγχες ότι είναι σωστός (σκέψου την αντιστροφή του);

\ex Αν η γωνία $\omega$ είναι οξεία, τι είδους είναι η $180^\circ-\omega$;

\zoni{Για το σπίτι}{ψάξε γωνίες}

\ex Βρες στο σπίτι πέντε αντικείμενα με ορθές γωνίες και τρία με αμβλείες.

\ex Φτιάξε με σπάγκο και κόμπους το τρίγωνο $3,\,4,\,5$ και έλεγξε με ορθογώνιο χαρτί αν η γωνία είναι ορθή.
